\par\medskip
\begin{center}\hypertarget{native-166-article119}{}\pdfbookmark[2]{119.}{article-119}119.\end{center}
Aus dem Zusammenhang zwischen den gegebenen und den unbekannten Größen unserer Aufgabe ließe sich diese leicht auf sechs Gleichungen oder sogar auf weniger zurückführen, da man die eine oder andere Unbekannte recht bequem eliminieren könnte. Dieser Zusammenhang ist jedoch so verwickelt, dass die Gleichungen äußerst schwer zu behandeln wären. Eine solche Trennung der Unbekannten, dass sich schließlich eine Gleichung mit nur einer Unbekannten ergibt, kann man im Allgemeinen\footnote{Liegen die Beobachtungen zeitlich so nahe beieinander, dass man die Zeitintervalle als unendlich kleine Größen behandeln darf, so gelingt eine solche Trennung, und die gesamte Aufgabe lässt sich auf die Lösung einer algebraischen Gleichung siebten oder achten Grades zurückführen.} für unmöglich halten. Umso weniger wird man die Aufgabe allein mit direkten Rechenoperationen vollständig lösen können.

Unsere Aufgabe lässt sich jedoch auf verschiedene Weise auf die Lösung von \emph{zwei} Gleichungen $X=0$, $Y=0$ zurückführen, in denen nur noch zwei Unbekannte $x$, $y$ miteinander verknüpft sind. Dabei müssen $x$, $y$ keine zwei der Elemente selbst sein. Sie können auf beliebige Weise mit den Elementen zusammenhängen, sofern sich die Elemente aus ihnen bequem ableiten lassen, sobald sie gefunden sind. Ebenso wenig ist es offenbar erforderlich, $X$, $Y$ als explizite Funktionen von $x$, $y$ darzustellen. Es genügt, dass sie durch ein Gleichungssystem so mit ihnen verbunden sind, dass man von gegebenen Werten von $x$, $y$ zu den entsprechenden Werten von $X$, $Y$ gelangen kann.
